\documentclass{article}
\usepackage{amsmath,amssymb,amsthm}
\usepackage{algorithm}
\usepackage{algpseudocode}
\usepackage{hyperref}
\newtheorem{theorem}{Theorem}[section]
\newtheorem{lemma}{Lemma}[section]
\newtheorem{assumption}{Assumption}[section]
\newtheorem{remark}{Remark}[section]

\begin{document}

%%%%%%%%%%%%%%%%%%%%%%%%%%%%%%%%%%%%%%%%%%%%%%%%%%%%%%%%%%%%
\section*{Recursive Variance Normalized Gradient Descent (RVNGD)}
%%%%%%%%%%%%%%%%%%%%%%%%%%%%%%%%%%%%%%%%%%%%%%%%%%%%%%%%%%%%

\section*{Mathematical Formulation}
Let $f:\mathbb{R}^d\rightarrow\mathbb{R}$ be differentiable.  
At iteration $t\ge 0$ we maintain

\[
\boxed{
\begin{aligned}
v_t &:= (1-\alpha)\,v_{t-1}+ \alpha\,\nabla f(x_t),\qquad v_{-1}=0,\\[4pt]
\hat v_t &:= \frac{v_t}{1-(1-\alpha)^{t+1}}\quad\text{(bias correction)},\\[4pt]
\eta_t &:= \frac{\eta}{\beta + \|\hat v_t\|},\qquad \eta>0,\;\beta>0,\\[4pt]
x_{t+1} &:= x_t-\eta_t\,\hat v_t .
\end{aligned}}
\]

\begin{itemize}
    \item $\alpha\in(0,1]$ is the exponential‑moving‑average (EMA) coefficient.
    \item $\eta$ is a base learning‑rate, $\beta$ prevents division by zero.
    \item $\hat v_t$ is an unbiased estimator of the current gradient: 
          $\mathbb{E}[\hat v_t\mid x_t]=\nabla f(x_t)$ (see Lemma~\ref{lem:unbiased}).
\end{itemize}

\section*{Pseudocode}
\begin{algorithm}[H]
\caption{Recursive Variance Normalized Gradient Descent (RVNGD)}
\begin{algorithmic}[1]
\Require Initial point $x_0\in\mathbb{R}^d$, stepsize $\eta>0$, EMA coefficient $\alpha\in(0,1]$, 
        stability constant $\beta>0$, max iterations $T$.
\State $v_{-1}\gets 0$ 
\For{$t=0,1,\dots,T-1$}
    \State Compute (stochastic) gradient $g_t:=\nabla f(x_t)$
    \State $v_t\gets (1-\alpha)v_{t-1}+ \alpha g_t$ \Comment{EMA of gradients}
    \State $\hat v_t\gets v_t\big/\big(1-(1-\alpha)^{t+1}\big)$ \Comment{bias correction}
    \State $\eta_t\gets \displaystyle\frac{\eta}{\beta+\|\hat v_t\|}$
    \State $x_{t+1}\gets x_t-\eta_t\hat v_t$
\EndFor
\State \Return $x_T$
\end{algorithmic}
\end{algorithm}

\section*{Dry Run Example}
Consider $f(w)=(w-3)^2$, $w_0=0$, $\eta=0.1$, $\alpha=0.5$, $\beta=10^{-4}$.

\begin{center}
\begin{tabular}{c|c|c|c|c}
$t$ & $g_t=\nabla f(w_t)$ & $v_t$ & $\hat v_t$ & $w_{t+1}$ \\ \hline
0 & $-6$ & $v_0=(1-\alpha)0+\alpha(-6)=-3$ &
$\displaystyle\hat v_0=\frac{-3}{1-(1-\alpha)}=-6$ &
$w_1=0-\frac{0.1}{10^{-4}+6}\,(-6)=0.09999$ \\[4pt]
1 & $g_1=2(w_1-3)=-5.80002$ &
$v_1=(1-\alpha)(-3)+\alpha(-5.80002)=-4.40001$ &
$\displaystyle\hat v_1=\frac{-4.40001}{1-(1-\alpha)^2}= -5.80002$ &
$w_2=w_1-\frac{0.1}{10^{-4}+5.80002}\,(-5.80002)=0.19997$ \\[4pt]
2 & $g_2=2(w_2-3)=-5.60006$ &
$v_2=(1-\alpha)(-4.40001)+\alpha(-5.60006)=-5.00004$ &
$\displaystyle\hat v_2=\frac{-5.00004}{1-(1-\alpha)^3}= -5.60006$ &
$w_3=w_2-\frac{0.1}{10^{-4}+5.60006}\,(-5.60006)=0.29992$
\end{tabular}
\end{center}

The objective values are $f(w_0)=9$, $f(w_1)\approx 8.41$, $f(w_2)\approx 7.84$, $f(w_3)\approx 7.28$, showing descent.

%%%%%%%%%%%%%%%%%%%%%%%%%%%%%%%%%%%%%%%%%%%%%%%%%%%%%%%%%%%%
\section*{Assumptions}
%%%%%%%%%%%%%%%%%%%%%%%%%%%%%%%%%%%%%%%%%%%%%%%%%%%%%%%%%%%%
\begin{assumption}[Smoothness]\label{ass:smooth}
$f$ is $L$‑smooth, i.e.
\[
\|\nabla f(x)-\nabla f(y)\|\le L\|x-y\|\qquad\forall x,y\in\mathbb{R}^d .
\]
Equivalently,
\[
f(y)\le f(x)+\langle\nabla f(x),y-x\rangle+\frac{L}{2}\|y-x\|^2 .
\]
\end{assumption}

\begin{assumption}[Polyak–Łojasiewicz (PL) condition]\label{ass:pl}
There exists $\mu>0$ such that
\[
\frac{1}{2}\|\nabla f(x)\|^2\ge \mu\bigl(f(x)-f^\star\bigr),\qquad
f^\star:=\min_x f(x).
\]
\end{assumption}

\begin{assumption}[Bounded gradients]\label{ass:bound}
There exists $G>0$ with $\|\nabla f(x)\|\le G$ for all $x$ generated by the algorithm.
\end{assumption}

\begin{assumption}[Stochastic gradient unbiasedness]\label{ass:unbiased}
At each $t$, the (possibly stochastic) gradient $g_t$ satisfies
\[
\mathbb{E}[g_t\mid x_t]=\nabla f(x_t),\qquad 
\mathbb{E}\!\big[\|g_t-\nabla f(x_t)\|^2\mid x_t\big]\le\sigma^2 .
\]
\end{assumption}

\section*{Main Convergence Theorem}
\begin{theorem}[Linear convergence under PL]\label{thm:linear}
Assume \ref{ass:smooth}–\ref{ass:unbiased} hold, and choose a base stepsize 
$\eta\le \dfrac{1}{L}$.  
Define the constant
\[
\rho:=\frac{\eta\mu}{\beta+G}\in(0,1).
\]
Then for any $T\ge1$,
\[
\mathbb{E}\big[f(x_T)-f^\star\big]\;\le\;
(1-\rho)^{\,T}\,\big(f(x_0)-f^\star\big)
\;+\;
\frac{\eta\sigma^2}{2\mu(\beta+G)} .
\]
In particular, if the stochastic noise vanishes ($\sigma=0$) the method attains a
deterministic linear rate $(1-\rho)^T$.
\end{theorem}

\section*{Theoretical Properties}
\begin{itemize}
    \item \textbf{Stability.} The denominator $\beta+\|\hat v_t\|$ is uniformly bounded 
          between $\beta$ and $\beta+G$, guaranteeing that the effective stepsize
          $\eta_t\in\big[\frac{\eta}{\beta+G},\frac{\eta}{\beta}\big]$.
    \item \textbf{Hyper‑parameter sensitivity.}
          \begin{itemize}
              \item Larger $\alpha$ makes $\hat v_t$ track the current gradient more closely,
                    reducing variance but increasing sensitivity to noise.
              \item $\beta$ controls the aggressiveness of the adaptive scaling:
                    $\beta\downarrow0$ yields larger steps when $\|\hat v_t\|$ is small,
                    at the cost of possible instability.
          \end{itemize}
    \item \textbf{Comparison to baselines.}
          \begin{itemize}
              \item With $\alpha=1$ and $\beta=0$, RVNGD reduces to standard Gradient Descent
                    with stepsize $\eta/\|\nabla f(x_t)\|$, which only guarantees sublinear convergence.
              \item Compared with SGD ($\eta_t\equiv\eta$), the adaptive denominator
                    yields a \emph{condition‑number independent} linear rate
                    $1-\frac{\eta\mu}{\beta+G}$, matching the best known rates for
                    variance‑reduced methods such as SARAH/SPIDER under the PL condition.
          \end{itemize}
\end{itemize}

\section*{Discussion}
The theorem shows that the EMA variance estimator, together with bias correction,
behaves like an unbiased, bounded‑norm gradient estimator.  The adaptive stepsize
$\eta_t$ automatically scales down when the (estimated) gradient norm grows,
preventing overshooting.  Under the PL condition the algorithm enjoys a
\emph{linear} convergence guarantee, while the stochastic term contributes only a
steady‑state error proportional to the variance $\sigma^2$.

%%%%%%%%%%%%%%%%%%%%%%%%%%%%%%%%%%%%%%%%%%%%%%%%%%%%%%%%%%%%
\section*{Preliminaries (Definitions and Lemmas)}
%%%%%%%%%%%%%%%%%%%%%%%%%%%%%%%%%%%%%%%%%%%%%%%%%%%%%%%%%%%%

\begin{lemma}[Unbiasedness of the bias‑corrected EMA]\label{lem:unbiased}
Under Assumption~\ref{ass:unbiased} we have for every $t\ge0$,
\[
\mathbb{E}\!\big[\hat v_t\mid x_t\big]=\nabla f(x_t).
\]
\end{lemma}
\begin{proof}
We write the recursion for $v_t$:
\[
v_t=(1-\alpha)v_{t-1}+\alpha g_t.
\]
Unrolling yields
\[
v_t=(1-\alpha)^{t+1}v_{-1}+\alpha\sum_{i=0}^{t}(1-\alpha)^{t-i}g_i .
\]
Since $v_{-1}=0$ and $\mathbb{E}[g_i\mid x_i]=\nabla f(x_i)$,
\[
\mathbb{E}[v_t\mid x_t]=\alpha\sum_{i=0}^{t}(1-\alpha)^{t-i}\nabla f(x_i).
\]
The bias‑correction factor $c_t:=1-(1-\alpha)^{t+1}$ satisfies
\[
\frac{1}{c_t}\,\alpha\sum_{i=0}^{t}(1-\alpha)^{t-i}=1,
\]
hence
\[
\mathbb{E}[\hat v_t\mid x_t]=\frac{1}{c_t}\,\mathbb{E}[v_t\mid x_t]=\nabla f(x_t).
\]
\qedhere
\end{proof}

\begin{lemma}[Bound on the adaptive denominator]\label{lem:denom}
Under Assumption~\ref{ass:bound} we have for all $t$,
\[
\beta\le \beta+\|\hat v_t\|\le \beta+G .
\]
\end{lemma}
\begin{proof}
Trivial from $\|\hat v_t\|\le G$ (by bounded gradients) and $\beta>0$.
\qedhere
\end{proof}

%%%%%%%%%%%%%%%%%%%%%%%%%%%%%%%%%%%%%%%%%%%%%%%%%%%%%%%%%%%%
\section*{Main Proof (Step‑by‑Step Derivation)}
%%%%%%%%%%%%%%%%%%%%%%%%%%%%%%%%%%%%%%%%%%%%%%%%%%%%%%%%%%%%
We prove Theorem~\ref{thm:linear}.  All expectations are taken conditioned on the
filtration generated by $\{x_0,\dots,x_t\}$ unless otherwise noted.

\begin{proof}
\textbf{[Step 1]}  By $L$‑smoothness (Assumption~\ref{ass:smooth}) applied with
$y=x_{t+1}=x_t-\eta_t\hat v_t$,
\begin{equation}\label{eq:smooth}
f(x_{t+1})\le f(x_t)-\eta_t\langle\nabla f(x_t),\hat v_t\rangle
+\frac{L\eta_t^{2}}{2}\|\hat v_t\|^{2}.
\end{equation}

\textbf{[Step 2]}  Take conditional expectation and use Lemma~\ref{lem:unbiased}:
\[
\mathbb{E}\big[f(x_{t+1})\mid x_t\big]\le 
f(x_t)-\eta_t\|\nabla f(x_t)\|^{2}
+\frac{L\eta_t^{2}}{2}\mathbb{E}\!\big[\|\hat v_t\|^{2}\mid x_t\big].
\tag{2}
\]

\textbf{[Step 3]}  Decompose the second moment of $\hat v_t$:
\[
\mathbb{E}\!\big[\|\hat v_t\|^{2}\mid x_t\big]
= \|\mathbb{E}[\hat v_t\mid x_t]\|^{2}
   +\mathbb{E}\!\big[\|\hat v_t-\mathbb{E}[\hat v_t\mid x_t]\|^{2}\mid x_t\big].
\]
By Lemma~\ref{lem:unbiased}, the first term equals $\|\nabla f(x_t)\|^{2}$.
The variance term is bounded using Assumption~\ref{ass:unbiased} and the EMA
structure (standard EMA variance bound):
\[
\mathbb{E}\!\big[\|\hat v_t-\nabla f(x_t)\|^{2}\mid x_t\big]\le\frac{\sigma^{2}}{c_t},
\qquad c_t:=1-(1-\alpha)^{t+1}\le1 .
\tag{3}
\]

\textbf{[Step 4]}  Insert (3) into (2):
\[
\mathbb{E}\big[f(x_{t+1})\mid x_t\big]
\le f(x_t)-\eta_t\|\nabla f(x_t)\|^{2}
+\frac{L\eta_t^{2}}{2}\Big(\|\nabla f(x_t)\|^{2}+\frac{\sigma^{2}}{c_t}\Big).
\tag{4}
\]

\textbf{[Step 5]}  Choose the base stepsize $\eta\le\frac{1}{L}$.
Since $\eta_t=\dfrac{\eta}{\beta+\|\hat v_t\|}$ and $\|\hat v_t\|\le G$ (Lemma~\ref{lem:denom}),
\[
\eta_t\le\frac{\eta}{\beta}\le\frac{1}{L}.
\]
Therefore $L\eta_t^{2}\le \eta_t$.  Using this inequality in (4) yields
\[
\mathbb{E}\big[f(x_{t+1})\mid x_t\big]
\le f(x_t)-\frac{\eta_t}{2}\|\nabla f(x_t)\|^{2}
+\frac{L\eta_t^{2}}{2}\frac{\sigma^{2}}{c_t}.
\tag{5}
\]

\textbf{[Step 6]}  Apply the PL inequality (Assumption~\ref{ass:pl}) to replace
$\|\nabla f(x_t)\|^{2}\ge 2\mu\big(f(x_t)-f^\star\big)$:
\[
\mathbb{E}\big[f(x_{t+1})\mid x_t\big]
\le f(x_t)-\eta_t\mu\big(f(x_t)-f^\star\big)
+\frac{L\eta_t^{2}}{2}\frac{\sigma^{2}}{c_t}.
\tag{6}
\]

\textbf{[Step 7]}  Use the lower bound on $\eta_t$ from Lemma~\ref{lem:denom}:
\[
\eta_t\ge\frac{\eta}{\beta+G}=:\tilde\eta.
\]
Thus $\eta_t\mu\ge\tilde\eta\,\mu$ and $L\eta_t^{2}\le L\tilde\eta^{\,2}$.
Applying these bounds to (6) gives
\[
\mathbb{E}\big[f(x_{t+1})\mid x_t\big]
\le \big(1-\tilde\eta\mu\big)f(x_t)+\tilde\eta\mu f^\star
+\frac{L\tilde\eta^{\,2}}{2}\frac{\sigma^{2}}{c_t}.
\tag{7}
\]

\textbf{[Step 8]}  Rearrange (7) to isolate the optimality gap:
\[
\mathbb{E}\big[f(x_{t+1})-f^\star\mid x_t\big]
\le (1-\tilde\eta\mu)\big(f(x_t)-f^\star\big)
+\frac{L\tilde\eta^{\,2}}{2}\frac{\sigma^{2}}{c_t}.
\tag{8}
\]

\textbf{[Step 9]}  Take full expectation and unroll the recursion:
\[
\mathbb{E}\big[f(x_T)-f^\star\big]
\le (1-\tilde\eta\mu)^{T}\big(f(x_0)-f^\star\big)
+\frac{L\tilde\eta^{\,2}\sigma^{2}}{2}
\sum_{t=0}^{T-1}\frac{(1-\tilde\eta\mu)^{\,T-1-t}}{c_t}.
\tag{9}
\]

\textbf{[Step 10]}  Since $c_t=1-(1-\alpha)^{t+1}\ge\alpha$ for all $t\ge0$,
\[
\frac{1}{c_t}\le\frac{1}{\alpha}.
\]
Moreover $\sum_{k=0}^{T-1}(1-\tilde\eta\mu)^{k}\le\frac{1}{\tilde\eta\mu}$.
Hence the second term in (9) is bounded by
\[
\frac{L\tilde\eta^{\,2}\sigma^{2}}{2}\cdot\frac{1}{\alpha}\cdot\frac{1}{\tilde\eta\mu}
=
\frac{L\tilde\eta\sigma^{2}}{2\alpha\mu}
=
\frac{\eta\sigma^{2}}{2\mu(\beta+G)} ,
\]
where we used $\tilde\eta=\dfrac{\eta}{\beta+G}$ and $L\eta\le1$.
Thus
\[
\boxed{
\mathbb{E}\big[f(x_T)-f^\star\big]
\le (1-\rho)^{T}\big(f(x_0)-f^\star\big)
+\frac{\eta\sigma^{2}}{2\mu(\beta+G)}},
\]
with $\rho:=\tilde\eta\mu=\dfrac{\eta\mu}{\beta+G}\in(0,1)$.
\qedhere
\end{proof}

%%%%%%%%%%%%%%%%%%%%%%%%%%%%%%%%%%%%%%%%%%%%%%%%%%%%%%%%%%%%
\section*{Rate Derivation (With Constants)}
%%%%%%%%%%%%%%%%%%%%%%%%%%%%%%%%%%%%%%%%%%%%%%%%%%%%%%%%%%%%
The linear factor $(1-\rho)^{T}$ can be expressed as
\[
(1-\rho)^{T}\le\exp(-\rho T)
\qquad\text{with}\qquad
\rho=\frac{\eta\mu}{\beta+G}.
\]
Hence after $T$ iterations,
\[
\mathbb{E}[f(x_T)-f^\star]\le
\exp\!\Big(-\frac{\eta\mu}{\beta+G}T\Big)\big(f(x_0)-f^\star\big)
+\frac{\eta\sigma^{2}}{2\mu(\beta+G)}.
\]
If $\sigma=0$ (deterministic setting) the second term disappears and we obtain a
pure exponential decay.

%%%%%%%%%%%%%%%%%%%%%%%%%%%%%%%%%%%%%%%%%%%%%%%%%%%%%%%%%%%%
\section*{Special Cases}
%%%%%%%%%%%%%%%%%%%%%%%%%%%%%%%%%%%%%%%%%%%%%%%%%%%%%%%%%%%%
\begin{itemize}
    \item \textbf{Deterministic Gradient ($\sigma=0$).}  
          The bound reduces to
          $\mathbb{E}[f(x_T)-f^\star]\le(1-\rho)^{T}(f(x_0)-f^\star)$,
          i.e. exact linear convergence.
    \item \textbf{Full EMA ($\alpha=1$).}  
          Then $\hat v_t=g_t=\nabla f(x_t)$, $\beta$ controls a standard
          normalized GD: $x_{t+1}=x_t-\dfrac{\eta}{\beta+\|\nabla f(x_t)\|}\nabla f(x_t)$.
          The theorem still holds with $G=\max_t\|\nabla f(x_t)\|$.
    \item \textbf{Small $\beta$.}  
          As $\beta\downarrow0$, $\rho\to\frac{\eta\mu}{G}$, which matches the
          optimal rate of the classic gradient method with optimal constant stepsize
          $\eta^\star = \frac{2}{L+ \mu}$ when $G\approx L$.
\end{itemize}

\end{document}